\documentclass%
%% versión borrador (o corrector) -> descomentar la siguiente linea
%[borrador]%
{tesisCNA}
%% Si no se quiere dedicar la tesis OPCIONAL -> descomentar la siguiente linea
%\def\nodedicatoria{1}
%% Si se va a hacer un índice alfabético (hay que ejecutar makeindex antes de lualatex !!), comentar la siguiente linea
\def\noindicealfabetico{1}
%% Si se quiere que aparezca la lista de figuras/tablas -> descomentar la siguiente linea
%\def\listatablasfiguras{1}%
\begin{document}
%% epígrafe OPCIONAL -> si no se desea, comentar la siguiente linea
\epigrafe
%%%% ¡¡¡ NO TOCAR LA SIGUIENTE LINEA !!!
\materialfijoinicial
% prólogo / preface OPCIONAL -> si no se desea, comentar la siguiente linea
\prefacio
% abreviaturas y símbolos / glossary OPCIONAL -> si no se desea, comentar la siguiente linea
\glosario
%%%% ¡¡¡ NO TOCAR LA SIGUIENTE LINEA !!!
\pagenumbering{arabic}%
%
%%%% Recordar que los objetivos deben presentarse explícitamente de forma obligatoria, como una sección en la introducción o como un capítulo aparte !!!!
\include{introduccion}
\cleardoublepage
%%
%%%% Si queremos un epígrafe antes del cap. 2, por ejemplo, descomentamos la siguiente linea
%\chapterepigraph{Je pense, donc je suis}{Discours de la méthode}{R.\ Descartes}
%%
\include{cap2}
\cleardoublepage
\include{cap3}
\cleardoublepage
%%%% Añádanse tantos capítulos como uno necesite
%
%%%% Recordar que en el formato de tesis por compendio, debe haber un capítulo de discusión global del trabajo que se incluye en los artículos antes de las conclusiones  !!!!
%
%\include{cap3}
%\cleardoublepage
%
\conclusionesperspectivas
%%% SI hay algún apéndice (sino quitar las tres lineas siguientes) 
\appendix
\include{apendice1}
\cleardoublepage
%\include{apendice2}
%\cleardoublepage
%%% terminan apéndices

%%%% BIBLIOGRAFÍA (se puede cambiar el estilo ... no se da soporte en ese caso...)
\bibliographystyle{cna}
\phantomsection
\addcontentsline{toc}{chapter}{\bibname}%
%% La bibliografía está en Plantilla.bib. Se recomienda usar un gestor de bibliografía ( p.ej. JABREF )
\bibliography{Plantilla}
\cleardoublepage 

%%% ÍNDICE ALFABÉTICO
\ifthenelse{\equal{\noindicealfabetico}{\null}}{\phantomsection\addcontentsline{toc}{chapter}{\indexname}\input Plantilla.ind\relax}{\relax}%
%
%%% en casos especiales (solo a veces es necesario en español) hay que usar esto. ¡¡¡ Preguntar !!!
%\input Plantilla.inx
%
%
% RESÚMENES (obligatorios). Pueden ponerse después de \materialfijoinicial
%  Uno tiene que ser la traducción del otro
%
\resumeneningles % MÁXIMO UNA PAGINA !! (el fichero es summary.tex)
\resumenenespanyol % MÁXIMO UNA PAGINA !! (el fichero es resumen.tex)
%
%%%% Y SE ACABO -- ¡¡¡ NO AÑADIR NADA MAS !!!
%
\end{document}
